% GENERATED FILE. Edit canonical lessons, then run npm run generate:books.
% canonical-source-hash: fnv1a64:5c1061833b21ca1b
% canonical-lessons: HI-C386-khisakna, HI-C386-khalna, HI-C386-sahara, HI-C386-zid, HI-C386-yogya

\chapter{Leaning on Others}
\label{ch:hi-leaning-on-others}
\hlchaptermodality{386}

\begin{chapteropening}
\textbf{By the end of this chapter:} \emph{I can say what will happen if someone leaves, stays or insists.}

Complete the last lesson of chapter 386: I can say what will happen if someone leaves, stays or insists.
\end{chapteropening}

\section[\texorpdfstring{\hi{खिसकना} (khisaknā)}{खिसकना (khisaknā)}]{\hi{खिसकना} (khisaknā) — to slip away, to shift}

\label{lesson:HI-C386-khisakna}



\begin{quote}
Before the new one: say the Hindi for to set, then the Hindi for unfinished.
\end{quote}

\subsection*{You'll want to know: \hi{खिसकना}}
\textbf{\hi{खिसकना}} — \emph{khisaknā} — \textquotedblleft{}to slip away, to shift\textquotedblright{}.

\begin{quote}
\hi{अगर} \hi{तुम} \hi{दोनों} \hi{थोड़ा} \hi{खिसकोगे} \hi{तो} \hi{मैं} \hi{भी} \hi{बैठ} \hi{सकूँगा।}
\end{quote}

\begin{quote}
\emph{agar tum donõ thoṛā khiskoge to maiṁ bhī baiṭh sakūṅgā.}
\end{quote}

\begin{quote}
If you two shift up a little, I'll be able to sit too.
\end{quote}

Past (he …): \hi{वह} \hi{पीछे} \hi{खिसक} \hi{गया}, \emph{vah pīche khisak gayā}.

\begin{cousinweb}
\textbf{\hi{खिसकना}} (\emph{khisaknā}) is inherited from \emph{*khis}, \textquotedblleft{}slip\textquotedblright{}, a form CDIAL reconstructs from its descendants (CDIAL 3888).
\end{cousinweb}

\begin{grammarlens}[title={What you've built}]
One more word for saying what happens if.
\end{grammarlens}

\subsection*{Your turn}
\begin{itemize}
\raggedright
  \item \emph{Say it:} \emph{khisaknā}
  \item \emph{Say it:} \emph{khisaknā}, once more
  \item \emph{Say it:} \emph{agar tum donõ thoṛā khiskoge to maiṁ bhī baiṭh sakūṅgā.} — If you two shift up a little, I'll be able to sit too.
  \item \emph{Recall:} say the Hindi for to set, then the Hindi for unfinished, then say \emph{khisaknā} again
\end{itemize}

\begin{tcolorbox}[breakable,colback=teal!4,colframe=teal!35!black,title={Before you move on}]
What does \hi{खिसकना} mean? (\textquotedblleft{}To slip away, to shift\textquotedblright{}.) Say it once more.
\end{tcolorbox}
\section[\texorpdfstring{\hi{खलना} (khalnā)}{खलना (khalnā)}]{\hi{खलना} (khalnā) — to be missed, to bother (someone)}

\label{lesson:HI-C386-khalna}



\begin{quote}
Before the new one: say the Hindi for unfinished, then the Hindi for to slip away.
\end{quote}

\subsection*{You'll want to know: \hi{खलना}}
\textbf{\hi{खलना}} — \emph{khalnā} — \textquotedblleft{}to be missed, to bother (someone)\textquotedblright{}.

\begin{quote}
\hi{अगर} \hi{तुम} \hi{चले} \hi{जाओगे} \hi{तो} \hi{मुझे} \hi{बहुत} \hi{खलेगा।}
\end{quote}

\begin{quote}
\emph{agar tum cale jāoge to mujhe bahut khalegā.}
\end{quote}

\begin{quote}
If you go away, I'll really miss you.
\end{quote}

Past (it …): \hi{वह} \hi{बात} \hi{उसे} \hi{बहुत} \hi{खली}, \emph{vah bāt use bahut khalī}.

\begin{cousinweb}
\textbf{\hi{खलना}} (\emph{khalnā}) is a tadbhava, inherited from Sanskrit \emph{skhalati}, \textquotedblleft{}stumbles\textquotedblright{} (CDIAL 13663). CDIAL glosses the Hindi verb \textquotedblleft{}to be bad, to be inimical\textquotedblright{}; the sense of being missed is not among its glosses.
\end{cousinweb}

\begin{grammarlens}[title={What you've built}]
One more word for saying what happens if.
\end{grammarlens}

\subsection*{Your turn}
\begin{itemize}
\raggedright
  \item \emph{Say it:} \emph{khalnā}
  \item \emph{Say it:} \emph{khalnā}, once more
  \item \emph{Say it:} \emph{agar tum cale jāoge to mujhe bahut khalegā.} — If you go away, I'll really miss you.
  \item \emph{Recall:} say the Hindi for unfinished, then the Hindi for to slip away, then say \emph{khalnā} again
\end{itemize}

\begin{tcolorbox}[breakable,colback=teal!4,colframe=teal!35!black,title={Before you move on}]
What does \hi{खलना} mean? (\textquotedblleft{}To be missed, to bother (someone)\textquotedblright{}.) Say it once more.
\end{tcolorbox}
\section[\texorpdfstring{\hi{सहारा} (sahārā)}{सहारा (sahārā)}]{\hi{सहारा} (sahārā) — a support, someone to lean on}

\label{lesson:HI-C386-sahara}



\begin{quote}
Before the new one: say the Hindi for to slip away, then the Hindi for to be missed.
\end{quote}

\subsection*{You'll want to know: \hi{सहारा}}
\textbf{\hi{सहारा}} — \emph{sahārā} — \textquotedblleft{}a support, someone to lean on\textquotedblright{}.

\begin{quote}
\hi{अगर} \hi{बूढ़ी} \hi{माँ} \hi{को} \hi{सहारा} \hi{नहीं} \hi{मिलेगा} \hi{तो} \hi{वह} \hi{गिर} \hi{जाएगी।}
\end{quote}

\begin{quote}
\emph{agar būṛhī māṁ ko sahārā nahīṃ milegā to vah gir jāegī.}
\end{quote}

\begin{quote}
If the old mother has nobody to lean on, she will fall.
\end{quote}

Masculine.

\begin{cousinweb}
\textbf{\hi{सहारा}} (\emph{sahārā}) is a tadbhava, inherited from Sanskrit \emph{sahakāra}, \textquotedblleft{}assistance\textquotedblright{} (CDIAL 13299).
\end{cousinweb}

\begin{grammarlens}[title={What you've built}]
One more word for saying what happens if.
\end{grammarlens}

\subsection*{Your turn}
\begin{itemize}
\raggedright
  \item \emph{Say it:} \emph{sahārā}
  \item \emph{Say it:} \emph{sahārā}, once more
  \item \emph{Say it:} \emph{agar būṛhī māṁ ko sahārā nahīṃ milegā to vah gir jāegī.} — If the old mother has nobody to lean on, she will fall.
  \item \emph{Recall:} say the Hindi for to slip away, then the Hindi for to be missed, then say \emph{sahārā} again
\end{itemize}

\begin{tcolorbox}[breakable,colback=teal!4,colframe=teal!35!black,title={Before you move on}]
What does \hi{सहारा} mean? (\textquotedblleft{}A support, someone to lean on\textquotedblright{}.) Say it once more.
\end{tcolorbox}
\section[\texorpdfstring{\hi{ज़िद} (zid)}{ज़िद (zid)}]{\hi{ज़िद} (zid) — stubbornness, insistence}

\label{lesson:HI-C386-zid}



\begin{quote}
Before the new one: say the Hindi for to be missed, then the Hindi for someone to lean on.
\end{quote}

\subsection*{You'll want to know: \hi{ज़िद}}
\textbf{\hi{ज़िद}} — \emph{zid} — \textquotedblleft{}stubbornness, insistence\textquotedblright{}.

\begin{quote}
\hi{ज़िद} \hi{मत} \hi{करो}, \hi{वरना} \hi{माँ} \hi{को} \hi{ग़ुस्सा} \hi{आएगा।}
\end{quote}

\begin{quote}
\emph{zid mat karo, varnā māṁ ko ġussā āegā.}
\end{quote}

\begin{quote}
Don't be stubborn, or your mother will get angry.
\end{quote}

Feminine.

\begin{cousinweb}
\textbf{\hi{ज़िद}} (\emph{zid}) is usually traced to Arabic \emph{ḍidd}, \textquotedblleft{}the opposite\textquotedblright{}, which reached Hindi through Persian.
\end{cousinweb}

\begin{grammarlens}[title={What you've built}]
One more word for saying what happens if.
\end{grammarlens}

\subsection*{Your turn}
\begin{itemize}
\raggedright
  \item \emph{Say it:} \emph{zid}
  \item \emph{Say it:} \emph{zid}, once more
  \item \emph{Say it:} \emph{zid mat karo, varnā māṁ ko ġussā āegā.} — Don't be stubborn, or your mother will get angry.
  \item \emph{Recall:} say the Hindi for to be missed, then the Hindi for someone to lean on, then say \emph{zid} again
\end{itemize}

\begin{tcolorbox}[breakable,colback=teal!4,colframe=teal!35!black,title={Before you move on}]
What does \hi{ज़िद} mean? (\textquotedblleft{}Stubbornness, insistence\textquotedblright{}.) Say it once more.
\end{tcolorbox}
\section[\texorpdfstring{\hi{योग्य} (yogya)}{योग्य (yogya)}]{\hi{योग्य} (yogya) — eligible, qualified}

\label{lesson:HI-C386-yogya}



\begin{quote}
Before the new one: say the Hindi for someone to lean on, then the Hindi for stubbornness.
\end{quote}

\subsection*{You'll want to know: \hi{योग्य}}
\textbf{\hi{योग्य}} — \emph{yogya} — \textquotedblleft{}eligible, qualified\textquotedblright{}.

\begin{quote}
\hi{अगर} \hi{तुम्हारे} \hi{अच्छे} \hi{अंक} \hi{आएँगे} \hi{तो} \hi{तुम} \hi{नौकरी} \hi{के} \hi{योग्य} \hi{होगे।}
\end{quote}

\begin{quote}
\emph{agar tumhāre acche aṅk āeṅge to tum naukrī ke yogya hoge.}
\end{quote}

\begin{quote}
If you get good marks, you'll be eligible for the job.
\end{quote}

\begin{cousinweb}
\textbf{\hi{योग्य}} (\emph{yogya}) is a tatsama, borrowed from Sanskrit \emph{yogya}, \textquotedblleft{}fit, suitable\textquotedblright{} (Monier-Williams).
\end{cousinweb}

\begin{grammarlens}[title={What you've built}]
One more word for saying what happens if.
\end{grammarlens}

\subsection*{Your turn}
\begin{itemize}
\raggedright
  \item \emph{Say it:} \emph{yogya}
  \item \emph{Say it:} \emph{yogya}, once more
  \item \emph{Say it:} \emph{agar tumhāre acche aṅk āeṅge to tum naukrī ke yogya hoge.} — If you get good marks, you'll be eligible for the job.
  \item \emph{Recall:} say the Hindi for someone to lean on, then the Hindi for stubbornness, then say \emph{yogya} again
\end{itemize}

\begin{tcolorbox}[breakable,colback=teal!4,colframe=teal!35!black,title={Before you move on}]
What does \hi{योग्य} mean? (\textquotedblleft{}Eligible, qualified\textquotedblright{}.) Say it once more.
\end{tcolorbox}
